% Imporditud: tools/import_eksperimendid.py
% Allikas: Eesti füüsikaolümpiaadi eksperimendi ülesannete
%   kogu (Rael Kalda, 2023), ülesanne Ü168, lahendus L168.
\setAuthor{EFO žürii}
\setRound{lõppvoor}
\setYear{2009}
\setNumber{G E2}
\setDifficulty{5}
\setTopic{vedelike mehaanika}
\setVahendid{süstlasilinder mahuskaalaga (cm$^3$), anum tundmatu vedelikuga, mõõtejoonlaud, statiiv klambriga, tühi plastpudel (ruumala V = 0,538 liitrit) tihedalt suletud korgiga, milles olevasse auku saab tihedalt suruda süstla otsa}
\setPunktid{14}
\setAllikas{Eksperimendi kogu Ü168, lahendus lk 126}
\setLahendusseis{toores}

\prob{Tundmatu vedelik}
Määrata vedeliku tihedus. Hinnata mõõtemääramatust.
\textit{Märkus:} katse on väga tundlik temperatuuri kõikumiste suhtes. Kirjeldage, milliseid meetmeid seda arvestades rakendasite.

\hint

\solu
Ühendame kokku süstla ja pudeli. Asetame süstla otsa vedelikku pressinud eelnevalt pudelist välja mõni cm$^3$ õhku. Süstalt üles alla liigutades leiame koha, kus vedeliku tase süstla sees ühtib vedeliku tasemega anumas. Sellisel juhul on õhurõhk pudelis võrdne välise õhurõhuga. Kehtib

m pV = RT,

M

kus V on süstla ja pudeli ruumalade summa. Kui nüüd sukeldame süstla sügavamale vedelikku, siis vedeliku samba rõhk suurendab õhurõhku pudelis. Kehtib

(V $-$ $\Delta$V )(p + $\Delta$p) = m RT , M

seda juhul, kui T = const. Vedeliku samba rõhk

$\Delta$p = $\rho$vg$\Delta$h,

kus $\Delta$h on vedeliku tasemete erinevus süstlas ja anumas, mille saame mõõtasüstla skaalalt joonlauaga. $\Delta$V on süsteemi pudel ja süstal ruumala muut, milles saame lugeda süstla mahuskaalalt. Kombineerides eelnevat kolme valemit saame,

pV = (V $- \Delta$V )(p + $\Delta$p) pV = (V $- \Delta$V ) (p + $\rho$vg$\Delta$h)

p$\Delta$V = (V $- \Delta$V )$\rho$vg$\Delta$h V $- \Delta$V $\approx$ V, väikese $\Delta$V korral.

Vedeliku tiheduseks saame: p$\Delta$V

$\rho$v = V . g$\Delta$h

Katseseade on väga tundlik temperatuuri kõikumiste suhtes. Selle vältimiseks kin-

nitada pudel klambrivahele, sest klambri temperatuur on sama, mis toas oleva õhu

temperatuur. Käega pudelist hoidmise korral soojeneb pudelis olev õhk. Näiteks

kui temperatuuri muut on $\Delta$T = 1 K, siis püsival rõhul on ruumala muut (tempe-

ratuuril 290 K )

V + $\Delta$V T + $\Delta$T $\Rightarrow \Delta$V $\approx$ 1,85 cm$^3$. = V T
\probend
